\lecture{6.5}{Writing Custom SuperLearner Wrappers}
\label{lec:6-5}

\subsection{Wrapper Function Structure}

A Super Learner wrapper is an R function that tells \texttt{SuperLearner()} how to fit
a candidate algorithm and how to generate predictions during cross-validation. Built-in
wrappers such as \texttt{SL.glm} follow the same structure custom wrappers must
follow.

The wrapper must return a list with two named components:
\[
\texttt{pred},\qquad \texttt{fit}.
\]
The \texttt{pred} component contains predictions on \texttt{newX}. The \texttt{fit}
component stores the fitted object or any information needed for later prediction.

\subsection{Required Inputs}

A standard wrapper has the form:

\begin{implementationbox}
\begin{lstlisting}[language=R]
SL.myalgorithm <- function(Y, X, newX, family, obsWeights, ...) {
  # fit model using Y and X
  # predict on newX
  # return list(pred = pred, fit = fit)
}
\end{lstlisting}
\end{implementationbox}

The key inputs are:
\begin{itemize}
\item \texttt{Y}: training outcomes;
\item \texttt{X}: training predictors;
\item \texttt{newX}: validation or prediction predictors;
\item \texttt{family}: outcome family;
\item \texttt{obsWeights}: observation weights, if supplied;
\item \texttt{...}: additional arguments passed by the package.
\end{itemize}

The \texttt{...} argument is important. It allows the wrapper to ignore arguments it
does not use without failing.

\subsection{Returning \texttt{pred} and \texttt{fit}}

The \texttt{pred} vector must have length \texttt{nrow(newX)}. The \texttt{fit} object
can be any list, but it should contain everything needed for future calls to
\texttt{predict()}.

\begin{implementationbox}
\begin{lstlisting}[language=R]
out <- list(pred = pred, fit = fit)
return(out)
\end{lstlisting}
\end{implementationbox}

It is useful to assign a class to the \texttt{fit} object:

\begin{implementationbox}
\begin{lstlisting}[language=R]
class(out$fit) <- "SL.myalgorithm"
\end{lstlisting}
\end{implementationbox}

The class allows R to dispatch to a custom prediction method later.

\subsection{Writing Prediction Methods}

When \texttt{predict()} is called on a Super Learner object containing a custom wrapper,
R must know how to predict from the custom fitted object. If the class is
\texttt{"SL.myalgorithm"}, R looks for a function named
\texttt{predict.SL.myalgorithm()}.

\begin{implementationbox}
\begin{lstlisting}[language=R]
predict.SL.myalgorithm <- function(object, newdata, ...) {
  pred <- predict(object$fitted_model, newdata = newdata, type = "response")
  return(pred)
}
\end{lstlisting}
\end{implementationbox}

The exact prediction call depends on the object stored inside \texttt{fit}. If
\texttt{fit} stores a \texttt{glm}, use \texttt{predict.glm}. If it stores a random
forest, use the prediction method for that random forest object.

\subsection{Custom Quadratic GLM Wrapper}

The following wrapper fits a linear regression with quadratic terms in $W_1$ and
$W_2$:

\begin{implementationbox}
\begin{lstlisting}[language=R]
SL.quadglm <- function(Y, X, newX, family, obsWeights, ...) {
  X <- as.data.frame(X)
  newX <- as.data.frame(newX)

  fit_glm <- glm(Y ~ W1 + I(W1^2) + W2 + I(W2^2) + A,
                 data = X,
                 family = family,
                 weights = obsWeights)

  pred <- predict(fit_glm, newdata = newX, type = "response")

  fit <- list(fitted_model = fit_glm)
  class(fit) <- "SL.quadglm"

  list(pred = pred, fit = fit)
}

predict.SL.quadglm <- function(object, newdata, ...) {
  predict(object$fitted_model, newdata = as.data.frame(newdata),
          type = "response")
}
\end{lstlisting}
\end{implementationbox}

This wrapper assumes the columns \texttt{W1}, \texttt{W2}, and \texttt{A} exist in
\texttt{X}. If those names are absent, the wrapper will fail. Custom wrappers should
therefore be written with clear expectations about input variable names.

\subsection{Custom Propensity Score Wrapper}

For the simulated data, the true propensity score has a $W_1W_2$ interaction. A custom
wrapper can fit that logistic regression:

\begin{implementationbox}
\begin{lstlisting}[language=R]
SL.truepropensity <- function(Y, X, newX, family, obsWeights, ...) {
  X <- as.data.frame(X)
  newX <- as.data.frame(newX)

  fit_glm <- glm(Y ~ W1 + W2 + W1:W2,
                 data = X,
                 family = binomial(),
                 weights = obsWeights)

  pred <- predict(fit_glm, newdata = newX, type = "response")

  fit <- list(fitted_model = fit_glm)
  class(fit) <- "SL.truepropensity"

  list(pred = pred, fit = fit)
}

predict.SL.truepropensity <- function(object, newdata, ...) {
  predict(object$fitted_model, newdata = as.data.frame(newdata),
          type = "response")
}
\end{lstlisting}
\end{implementationbox}

This learner can be included in a propensity score library:

\begin{implementationbox}
\begin{lstlisting}[language=R]
sl_fit_g <- SuperLearner(
  Y = A,
  X = data.frame(W1 = W1, W2 = W2),
  SL.library = c("SL.glm", "SL.mean", "SL.truepropensity"),
  family = binomial(),
  method = "method.CC_nloglik",
  cvControl = list(V = 2)
)
\end{lstlisting}
\end{implementationbox}

\subsection{Logistic-Linear Ensembles}

For binary outcomes or treatments, \texttt{method.CC\_nloglik} uses a
logistic-linear ensemble. If $g_{n,k}(w)$ is the predicted probability from candidate
$k$, the ensemble has the form
\[
g_{n,\omega}(w)
=
\expit\left\{\sum_{k=1}^K\omega_k\logit(g_{n,k}(w))\right\}.
\]
This combines candidates on the log odds scale and maps the result back to a
probability. The weights still satisfy $\omega_k\ge 0$ and $\sum_k\omega_k=1$.

\begin{keyidea}
A wrapper is the contract between an algorithm and \texttt{SuperLearner()}: fit on
training data, predict on new data, and store enough information to predict again
later.
\end{keyidea}
